\subsection{The authentication service}
\label{sec:ha:disabled:auth}

Several authentication services run at once, and none of them leads. A member
never contacts one directly. A gateway contacts an authentication service on the
member's behalf at logon, and reaches whichever instance answers.

Losing one instance reduces the capacity available and does not affect
correctness. Two instances holding different answers to the same question does
affect correctness, and most of this subsection concerns that condition.

\subsubsection{What state it holds}
\label{sec:ha:disabled:auth:state}

An instance holds the credential set: which comp ids may trade, the material each
is authenticated against, and the terms each was admitted under. This set is a
copy rather than the original. The database holds the original, and the
administration service is the only component that writes to it. That service
changes credentials and distributes each change to every authentication service,
and \cref{sec:ha:disabled:admin} describes it.

\subsubsection{Where the copy that outlives it lives}
\label{sec:ha:disabled:auth:durable}

The database holds the copy, and an instance reaches it in two steps rather than
one. An instance loads an export of the database when it starts, and the
administration service afterwards sends it each change as it makes it. The
instance holds no connection to the database, so it can neither check an answer
against the original nor establish how far behind it has fallen.

\textbf{Only the administration service can establish whether an instance is
current}, because only that service holds a record of what it sent and to which
instance.

\begin{gap}{}
  A change is reported as having succeeded when at least one instance applied
  it. A partial fan-out is therefore indistinguishable from a complete one, and
  nobody learns that an instance was missed.
\end{gap}

\begin{gap}{}
  Nothing brings a running instance up to date. Having missed a change, it stays
  wrong until it is restarted --- and restarting reloads the export, which may
  itself predate the change, so restarting is not a reliable remedy either.
\end{gap}

\subsubsection{How a restarted instance becomes current}
\label{sec:ha:disabled:auth:recovery}

The instance reloads the export and resumes answering. It has one such operation
rather than two, as the gateway in \cref{sec:ha:disabled:gateway} does, because
an authentication service holds no role and is never promoted. A restarted
instance and a newly started one therefore do the same thing.

The instance does not ask for the changes it missed, and no component offers them
to it.

\subsubsection{What the member is told}
\label{sec:ha:disabled:auth:member}

An instance can be wrong in two ways, and the two are not equally serious, which
is why the requirements below treat them differently. An instance that has missed
an addition refuses a member that should be admitted. The member sees that
refusal at once, and the venue can recover from it. An instance that has missed a
withdrawal admits a member that should be refused. Nothing reports that
admission, and it persists until somebody corrects the instance.

\begin{requirement}{R-0042}{A credential change is applied by every instance, or is reported as incomplete}
  A change shall not be reported as having succeeded unless every instance
  applied it. Where an instance did not, the change shall be reported as
  incomplete and shall name the instance.
  \rationale{Reporting success when one instance of several applied the change leaves a
  partial distribution looking the same as a complete one, so nobody can act on it.
  Every remedy below depends on the venue knowing that the condition exists.}
  \uncovered{no scenario makes a credential change while an instance is down}
\end{requirement}

\begin{requirement}{R-0043}{An instance that has missed a change is brought current without being restarted}
  Where an instance has missed a change, it shall be brought up to date while
  running.
  \rationale{Restarting an instance is not a remedy, because the instance reloads an export of
  the database that may itself be older than the change it missed. Nothing else
  brings an instance up to date, so without this requirement an instance stays
  stale indefinitely.}
  \uncovered{no scenario makes an instance miss a change, and nothing would bring it current}
\end{requirement}

\begin{requirement}{R-0044}{A member is admitted by every instance or by none}
  A member's logon shall succeed or be refused alike whichever instance answers
  it.
  \rationale{Intermittent admission is worse for a member than consistent refusal. The logon
  succeeds, then it fails, and nothing distinguishes the two attempts, so the
  member can neither report the fault clearly nor reproduce it.}
  \uncovered{no scenario authenticates the same member against both instances}
\end{requirement}

\begin{requirement}{R-0045}{An instance can say how current its credential set is}
  An instance shall be able to report the point up to which it has applied
  changes.
  \rationale{Where an instance cannot report this, no component can establish it. The instance
  has no view of the database, and the administration service cannot distinguish an
  instance that missed a change from one that never received it.}
  \uncovered{no scenario asks an instance how current its credentials are, and nothing answers}
\end{requirement}

\begin{requirement}{R-0046}{A change made while an instance is unavailable is applied when it returns}
  A credential change shall be accepted while an instance is unavailable, that
  instance shall be recorded as not current, and the change shall be applied to
  it when it returns.
  \rationale{Refusing the change instead would stop the venue withdrawing access during an
  outage, which is when withdrawing access may matter most. The cost is an instance
  that is behind for a bounded period. That cost is acceptable because
  \reqref{R-0042} makes the venue aware of the condition.}
  \uncovered{no scenario makes a credential change while an instance is down}
\end{requirement}
